\documentclass[aspectratio=169,11pt]{beamer}
\usepackage[T1]{fontenc}
\usepackage[utf8]{inputenc}
\usepackage[provide=*,polish]{babel}
\usepackage{lmodern}
\usepackage{amsmath}
\usepackage{booktabs}
\usepackage{graphicx}
\usepackage{hyperref}

\definecolor{navy}{HTML}{152A43}
\definecolor{teal}{HTML}{007F86}
\definecolor{light}{HTML}{EFF5F6}
\setbeamercolor{normal text}{fg=navy,bg=white}
\setbeamercolor{frametitle}{fg=navy,bg=white}
\setbeamercolor{structure}{fg=teal}
\setbeamercolor{block title}{fg=white,bg=teal}
\setbeamercolor{block body}{fg=navy,bg=light}
\setbeamertemplate{navigation symbols}{}
\setbeamertemplate{footline}[frame number]
\setbeamertemplate{itemize items}[circle]
\setbeamersize{text margin left=11mm,text margin right=11mm}

\title{MuJoCo HalfCheetah}
\subtitle{Uczenie ze wzmocnieniem \quad | \quad Drugi punkt kontrolny}
\author{Aleksandra Wieczorek \and Wojciech W\'ojcik}
\institute{Politechnika Warszawska}
\date{19 listopada 2026}

\begin{document}

\begin{frame}
  \titlepage
  \vspace{-0.5em}
  \begin{center}\small Temat zielony\end{center}
\end{frame}

\begin{frame}{Co zrealizowaliśmy}
  \begin{columns}[T,onlytextwidth]
    \begin{column}{0.48\textwidth}
      \begin{block}{Własna implementacja SAC}
        \begin{itemize}
          \item Implementacja PyTorch w \texttt{src/sac.py}.
          \item Polityka Gaussa z \texttt{tanh} squashing.
          \item Twin critics i replay buffer.
          \item Miękkie aktualizacje sieci docelowej.
          \item Automatyczna entropia \texttt{ent\_coef=auto}.
        \end{itemize}
      \end{block}
    \end{column}
    \begin{column}{0.48\textwidth}
      \begin{block}{Gotowy potok eksperymentalny}
        \begin{itemize}
          \item Konfigurowalny trening, ewaluacja i zapis modeli.
          \item Automatyczne katalogi dla każdego seeda.
          \item Ewaluacja deterministyczna bez eksploracji.
          \item Krzywa uczenia ze średnią i odchyleniem standardowym.
        \end{itemize}
      \end{block}
    \end{column}
  \end{columns}
\end{frame}

\begin{frame}{Własna implementacja SAC}
  \begin{columns}[T,onlytextwidth]
    \begin{column}{0.55\textwidth}
      \begin{itemize}
        \item \textbf{Actor}: sieć $17 \rightarrow 256 \rightarrow 256 \rightarrow 6$.
        \item \textbf{Critics}: dwa modelowane estymatory $Q(s,a)$.
        \item Akcje są ograniczane do przedziału $[-1,1]$ przez \texttt{tanh}.
        \item Próbki trenowane są pobierane z bufora doświadczeń.
        \item Model zapisuje i odczytuje własny format \texttt{mujoco-half-sac-v1}.
      \end{itemize}
    \end{column}
    \begin{column}{0.41\textwidth}
      \begin{block}{Cel SAC}
        \centering
        \vspace{0.3em}
        $\pi^*=\arg\max_\pi
        \sum_t \mathbb{E}\left[r_t+\alpha\,\mathcal{H}(\pi(\cdot|s_t))\right]$
        \vspace{0.3em}
      \end{block}
      \vspace{0.5em}
      \small Entropia zachęca do eksploracji, a uczenie off-policy
      ponownie wykorzystuje doświadczenia z replay buffer.
    \end{column}
  \end{columns}
\end{frame}

\begin{frame}{Konfiguracja}
  \centering
  \begin{tabular}{@{}lrp{0.42\textwidth}@{}}
    \toprule
    Parametr & Wartość & Znaczenie \\
    \midrule
    \texttt{total\_timesteps} & 100\,000 & Budżet kroki treningu \\
    \texttt{learning\_starts} & 10\,000 & Początek aktualizacji sieci \\
    \texttt{buffer\_size} & 100\,000 & Rozmiar replay buffer \\
    \texttt{batch\_size} & 256 & Rozmiar mini-paczki \\
    \texttt{learning\_rate} & $3\times10^{-4}$ & Tempo uczenia \\
    \texttt{gamma} & 0{,}99 & Dyskontowanie nagród \\
    \texttt{tau} & 0{,}005 & Miękka aktualizacja target networks \\
    \texttt{seeds} & 47, 69, 123 & Powtórzenia eksperymentu \\
    \bottomrule
  \end{tabular}
  \vspace{0.7em}

\end{frame}

\begin{frame}{Metodologia ewaluacji}
  \begin{columns}[T,onlytextwidth]
    \begin{column}{0.48\textwidth}
      \begin{block}{Tryb ewaluacji}
        \begin{itemize}
          \item \texttt{predict(..., deterministic=True)},
          \item brak eksploracji i bez dalszego uczenia,
          \item pełne epizody,
          \item osobny seed początkowy ewaluacji,
          \item osobne pliki dla każdego modelu.
        \end{itemize}
      \end{block}
    \end{column}
    \begin{column}{0.48\textwidth}
      \begin{block}{Wybierane metryki}
        \begin{itemize}
          \item zwrot z epizodu,
          \item nagroda naprzód i koszt sterowania,
          \item prędkość i przemieszczenie w osi $x$,
          \item średni kwadrat działania,
          \item długość epizodu,
          \item zakończenie i limit czasu.
        \end{itemize}
      \end{block}
    \end{column}
  \end{columns}
\end{frame}

\begin{frame}{Krzywe uczenia i wykresy}
  \begin{itemize}
    \item Skrypt \texttt{make\_plots.py} pobiera pliki
    \texttt{training\_metrics.csv} dla wszystkich seedów z konfiguracji.
    \item Wyniki są grupowane według numeru epizodu.
    \item Wykres pokazuje średnią zwrotu pomiędzy seedami.
    \item Pasmo wokół krzywej przedstawia odchylenie standardowe.
    \item Wykres można zapisać jako PNG do katalogu \texttt{artifacts/plots/}.
  \end{itemize}

  \vspace{0.7em}
  \begin{block}{Cel analizy}
    Ocenienie trendu uczenia oraz zmienności pomiędzy różnymi powtórzeniami,
    bez testowania osobnych konfiguracji jako jednego przebiegu.
  \end{block}
\end{frame}

\begin{frame}{Sprawdzenie implementacji}
  \begin{columns}[T,onlytextwidth]
    \begin{column}{0.48\textwidth}
      \begin{block}{Testy}
        \begin{itemize}
          \item 7 testów jednostkowych i integracyjnych przechodzi.
          \item Testy środowiska i metryk.
          \item Test zapisu i odczytu własnego modelu SAC.
          \item Test próbkowania i aktualizacji sieci.
        \end{itemize}
      \end{block}
    \end{column}
    \begin{column}{0.48\textwidth}
      \begin{block}{Weryfikacja end-to-end}
        \begin{itemize}
          \item Krótkie treningi dla trzech seedów uruchomiły się poprawnie.
          \item Ewaluacja załadowała i uruchomiła model każdego seeda.
          \item Został wygenerowany i zapisany wykres krzywych.
        \end{itemize}
      \end{block}
    \end{column}
  \end{columns}
\end{frame}

\begin{frame}{Podsumowanie punktu kontrolnego}
  \begin{enumerate}
    \setlength{\itemsep}{0.7em}
    \item Zaimplementowano własnego SAC w PyTorch i zastąpiono nim bibliotekę zewnętrzną.
    \item Uporządkowano trening, ewaluację, checkpointy i metryki.
    \item Dodano konfigurację listy seedów i osobne katalogi dla każdego powtórzenia.
    \item Przygotowano deterministyczną metodologię ewaluacji i zestaw metryk.
    \item Dodano generowanie krzywych uczenia oraz wizualizację polityki.
  \end{enumerate}
\end{frame}

\end{document}
